\documentclass[12pt,a4paper]{article}

\usepackage[utf8]{inputenc}
\usepackage[T5]{fontenc}
\usepackage[vietnamese,english]{babel}
\usepackage{lmodern}
\usepackage[top=2.5cm, bottom=2.5cm, left=2.5cm, right=2.5cm]{geometry}
\usepackage{enumitem}
\usepackage{amssymb}
\usepackage{booktabs}
\usepackage{graphicx}
\usepackage{hyperref}
\usepackage{setspace}
\usepackage{xcolor}
\usepackage{tabularx}

\onehalfspacing
\pagestyle{empty}

\begin{document}

% ============================================================
% HEADER
% ============================================================
\noindent
\begin{minipage}[c]{0.15\textwidth}
    \includegraphics[width=\linewidth]{Logo.png}
\end{minipage}%
\hfill
\begin{minipage}[c]{0.82\textwidth}
    \begin{center}
        {\normalsize\textbf{UNIVERSITY OF SCIENCE \& TECHNOLOGY OF HANOI}} \\
        {\normalsize\textbf{TRƯỜNG ĐẠI HỌC KHOA HỌC \& CÔNG NGHỆ HÀ NỘI}}
    \end{center}
\end{minipage}

\vspace{0.8cm}
\begin{center}
    {\Large\textbf{INTERNSHIP TOPIC PROPOSAL}} \\[0.5cm]
\end{center}

\noindent
\textbf{Specialty:} \quad $\boxtimes$ Information and Communication Technology \hfill $\square$ Cyber Security \\[0.3cm]
\textbf{Academic year:} 2025--2026

\vspace{0.5cm}
\noindent\rule{\textwidth}{0.8pt}

% ============================================================
% STUDENT INFORMATION
% ============================================================
\section*{STUDENT INFORMATION}

\begin{tabularx}{\textwidth}{@{}lX}
    \textbf{Full name:} & Phạm Hoàng An \\[0.2cm]
    \textbf{Student ID:} & 23BI14002 \\
\end{tabularx}

\vspace{0.5cm}
\noindent\rule{\textwidth}{0.4pt}

% ============================================================
% TOPIC
% ============================================================
\section*{Topic}

Enhanced Few-Shot Open-Set Keyword Spotting with Noise-Robust Prototype Classification and Real-Time Streaming Inference

% ============================================================
% DESCRIPTION
% ============================================================
\section*{Description}

Keyword Spotting (KWS) is the task of detecting predefined spoken keywords in audio streams. It is the core technology behind voice-activated systems such as ``Hey Siri,'' ``OK Google,'' and smart home controllers. Traditional KWS systems require thousands of labeled audio samples per keyword and can only recognize a fixed set of words defined during training.

This project builds a KWS system that overcomes these limitations by combining two key capabilities:

\begin{itemize}[leftmargin=1.5cm]
    \item \textbf{Few-Shot Learning:} The system can learn to recognize a new keyword from only 3--5 audio samples, without retraining the neural network. This is achieved through Prototypical Networks: each keyword is represented by a ``prototype'' (the average embedding of its few samples), and new audio is classified by measuring its distance to all prototypes.

    \item \textbf{Open-Set Recognition:} The system can reject unknown words that do not belong to any enrolled keyword. Instead of forcing a classification among known classes, it uses an acceptance radius threshold around each prototype --- if the distance exceeds this threshold, the input is rejected as ``unknown.''
\end{itemize}

The main contribution is a \textbf{Direct L2 distance classification} method that replaces probability-based scoring, providing a more intuitive and effective open-set rejection mechanism where the threshold has a clear geometric meaning as an ``acceptance radius.''

The system pipeline consists of:
\begin{enumerate}[leftmargin=1.5cm]
    \item \textbf{Audio preprocessing:} Raw audio (16\,kHz, 1-second segments) is converted to MFCC features (40 coefficients computed, 10 retained after narrowing and transposition).
    \item \textbf{Feature encoding:} A lightweight Depthwise Separable CNN (DSCNN-L, 276 channels) maps MFCC features to a 276-dimensional embedding space, trained with Triplet Loss on the MSWC English dataset.
    \item \textbf{Classification:} An Open Nearest Class Mean (openNCM) classifier with Direct L2 distance determines keyword identity or rejects unknown inputs.
    \item \textbf{Noise robustness (Extension 1):} Noise augmentation during training (DEMAND dataset, fixed 5\,dB SNR) and optional audio denoising at inference time.
    \item \textbf{Streaming inference (Extension 2):} Voice Activity Detection (Silero VAD) combined with a sliding-window mechanism enables continuous real-time keyword detection.
\end{enumerate}


% ============================================================
% EXPECTED OUTCOMES
% ============================================================
\section*{Expected Outcomes}

\begin{enumerate}[leftmargin=1.5cm]
    \item A fully functional few-shot open-set KWS system capable of learning new keywords from 3--5 samples without retraining.

    \item Quantitative evaluation comparing the proposed Direct L2 distance method against a probability-based baseline, using standard metrics: DET curves, AUC, and ACC@5\%FAR, across multiple evaluation protocols (GSC Fixed, GSC Randomized, MSWC Randomized).

    \item A noise robustness benchmark showing KWS accuracy at different SNR levels (0, 5, 10, 20\,dB and clean), with and without denoising preprocessing.

    \item A real-time streaming inference module with end-to-end latency under 100\,ms, integrating VAD and sliding-window detection.

    \item Complete source code, trained model checkpoints, and a written thesis report.
\end{enumerate}


% ============================================================
% USED METHODS AND TECHNIQUES
% ============================================================
\section*{Used Methods and Techniques}

\begin{tabularx}{\textwidth}{@{}l X}
    \toprule
    \textbf{Category} & \textbf{Details} \\
    \midrule
    Feature extraction &
        MFCC (40 coefficients $\rightarrow$ narrow to 10, transpose);
        \texttt{torchaudio} library \\[0.3cm]
    Neural network &
        DSCNN-L (Depthwise Separable CNN, 276 channels, 5 DS blocks,
        276-dim embeddings); lightweight architecture suitable for
        on-device deployment \\[0.3cm]
    Training method &
        Episodic metric learning with Triplet Loss (margin=0.5);
        Adam optimizer with StepLR scheduler ($\gamma=0.5$, step=20);
        40 epochs $\times$ 400 episodes \\[0.3cm]
    Classification &
        Prototypical Networks with Direct L2 distance;
        Open Nearest Class Mean (openNCM) for open-set rejection;
        threshold-based acceptance radius \\[0.3cm]
    Noise robustness &
        Training-time noise augmentation (DEMAND dataset, $p=0.95$,
        SNR=5\,dB); inference-time spectral gating denoising \\[0.3cm]
    Streaming &
        Silero VAD for voice activity detection;
        sliding window (1\,s window, 0.5\,s stride) \\[0.3cm]
    Demo (optional) &
        Gradio web application for interactive keyword enrollment,
        offline detection, and streaming demonstration \\[0.3cm]
    Datasets &
        Google Speech Commands v2 (GSC, 35 words -- evaluation);
        Multilingual Spoken Words Corpus (MSWC English, 763 words -- training);
        DEMAND (noise augmentation) \\[0.3cm]
    Evaluation &
        DET curves, AUC, ACC@5\%FAR, FRR@5\%FAR;
        3 protocols (GSC Fixed, GSC Randomized, MSWC Randomized);
        5--10 runs for statistical reliability \\[0.3cm]
    Tools \& frameworks &
        Python 3.10+, PyTorch 2.0+, torchaudio, SpeechBrain,
        Gradio, matplotlib, TensorBoard \\[0.3cm]
    Development &
        GPU training on Google Colab / Kaggle;
        Cursor IDE with AI-assisted coding rules \\
    \bottomrule
\end{tabularx}

\vspace{1.5cm}

% ============================================================
% SIGNATURES
% ============================================================
\noindent
\begin{tabularx}{\textwidth}{XX}
    \centering\textbf{Student's Signature} & \centering\textbf{Supervisor's Signature} \tabularnewline
    & \tabularnewline
    & \tabularnewline
    & \tabularnewline
    \centering Phạm Hoàng An & \centering\arraybackslash \\
\end{tabularx}

\end{document}
